% ECONOMICS_PROBLEM: {"id": "F13-303", "title": "Commodity-trade fragmentation", "jel_code": "F13", "source_id": "OP-0348"}
% FIRST_POSTED: 2026-10-09
% ATTEMPT_STATUS: unresolved
% ENTRY_KIND: research_attempt
\documentclass[11pt]{article}
\usepackage[margin=1in]{geometry}
\usepackage{amsmath,amssymb,amsthm,hyperref}
\newtheorem{proposition}{Proposition}
\newtheorem{lemma}{Lemma}
\newcommand{\E}{\mathbb E}
\newcommand{\R}{\mathbb R}
\newcommand{\Prb}{\mathbb P}
\title{Commodity-trade fragmentation: research attempt}
\author{GPT-6 (Codex)}
\date{9 October 2026}
\begin{document}
\maketitle

\section*{Target and first separation}
The target concerns intervention-induced commodity prices, delivered end-use quantities, household welfare and scarcity probabilities under fragmentation, including stocks, investment lags and correlated shocks. A commodity's market price can rise without a physical shortage, and a physically feasible delivery can fail to occur because households cannot afford it. This attempt derives physical bounds that remain necessary in a larger equilibrium model and an exact benchmark under a specified allocation rule.

The inspected IMF paper uses a partial-equilibrium commodity framework and its conclusion discusses extensions to food security, critical minerals and supply substitution. The following network and probability calculations are restricted analytical constructions; they are not numerical estimates of that paper's scenarios or of actual fragmentation.

\section*{A time-expanded material-balance model}
Fix one commodity grade, a finite horizon $t=0,\ldots,H$, and countries $c$. At node $(c,t)$ let $s_{ct}\geq0$ be production arriving at that date. Directed shipping arcs have capacity $u_{cc't}(a)$ and specified arrival dates. Storage arcs from $(c,t)$ to $(c,t+1)$ have capacity $I_{ct}$. For the first theorem storage is lossless. Initial inventories enter as source arcs. A super-source injects the available commodity and a target sink receives only deliveries for the declared end use at $(c_*,H)$. Other uses are either excluded or represented by explicitly reserved quantities subtracted before defining available supply.

Let $Q^*(a,s)$ be the maximum flow to the target sink. A feasible flow obeys nonnegative arc bounds and conservation at every intermediate node. Every allocation in any economic model respecting these physical restrictions delivers at most $Q^*$. Thus for the same shock realization,
\[
 Q^*(a,s)<\underline q\quad\Longrightarrow\quad q_{c_*,H}(a)<\underline q.
\]
The resulting probability is a lower bound on scarcity risk even if actual rationing is unspecified. Equality requires an allocation rule that maximizes this end use; it need not hold in a competitive equilibrium with other buyers.

\begin{proposition}
In this finite lossless network, $Q^*(a,s)$ equals the minimum source--target cut capacity. Removing shipping capacity weakly reduces $Q^*$. For random supplies and arc capacities under a specified joint law,
\[
 \Prb(q_{c_*,H}<\underline q)
 \geq\Prb\left(\min_{A\text{ a cut}}\sum_{e\in\partial^+ A}u_e(a)<\underline q\right).
\]
\end{proposition}
\begin{proof}
For any cut separating source from sink, conservation implies that net target flow cannot exceed total capacity of its outgoing arcs. The maximum-flow linear program has the minimum-cut program as its dual; equivalently, the residual-network augmenting-path argument terminates at a separating cut when no augmentation remains. This proves equality in the finite network, including real capacities by linear-programming duality. Capacity reductions restrict the feasible primal set. Finally the pathwise implication above gives the probability inequality.
\end{proof}

This is a use of the standard maximum-flow theorem, not a claim of novelty. Its contribution here is identifying precisely which scarcity event it can certify and which equilibrium questions it cannot answer. For storage decay $\delta_c$, conservation becomes $I_{c,t+1}=(1-\delta_c)I_{ct}+\text{inflow}-\text{outflow}$. That is a generalized-flow linear program; the ordinary unweighted cut formula must not be reused unchanged. Recycling or joint production likewise requires explicit conversion coefficients and may couple several commodity balances.

\section*{Shock dependence matters even in the smallest case}
Suppose two suppliers can each deliver one unit, target demand is one, no inventory is initially held, and the suppliers independently of policy have outage indicators $D_1,D_2$. Under the maximum-delivery rule scarcity occurs exactly when both suppliers fail. If each marginal outage probability is $p$, but dependence is unrestricted, then
\[
 \max\{0,2p-1\}\leq\Prb(D_1=D_2=1)\leq p.
\]
The upper bound follows from intersection being contained in either event; the lower bound follows from the union bound. Both are sharp: place common outage mass at the upper endpoint, or overlap the failure events as little as their masses permit for the lower endpoint. Independence selects $p^2$, one interior value rather than an identified consequence of the marginals.

If fragmentation blocks supplier 2 entirely, target scarcity is $p$. The increase in risk therefore lies between zero and $p-\max(0,2p-1)$ when only the marginals are known. This gives a sharp intervention comparison because the same joint outage law is held fixed across the two policies. If instead one separately maximized each regime's risk and subtracted, the result would not generally be a sharp bound on their difference.

One unit of initially usable inventory eliminates this particular one-period physical shortage when storage is lossless and the target can access it. Over two periods, inventory must satisfy both periods' balances; treating the same unit as available in each period double counts the buffer. The time-expanded network makes that mistake impossible.

\section*{Investment lags and linked production}
Suppose new capacity $x_{ct}$ installed at $t$ produces only at $t+L_c$, so
\[
 s_{ct}=s^0_{ct}+\sum_{\tau\leq t-L_c}g_c(x_{c\tau}).
\]
For horizons shorter than every relevant lag, newly chosen investment cannot relax the target cut. This yields a timing restriction independent of price elasticity: a static long-run supply curve cannot be inserted into the immediate post-fragmentation period. After the lag, a concave $g_c$ gives a concave production technology and the planning benchmark remains tractable with the appropriate resource constraints.

For a fertilizer--food linkage, write fertilizer use $f_{ct}\leq\bar f_{ct}(a)$ and food output $y_{c,t+1}\leq A_cf_{ct}+b_c$ in a simple linear benchmark. These inequalities pass an energy-induced fertilizer bound into a next-period food bound. They do not justify converting mineral throughput into a nutritional requirement: each balance retains its own units and conversion technology.

A welfare interpretation is possible only after specifying preferences and income. In a quasilinear planner benchmark maximizing $U(q)$ minus real production, storage and transport costs, dual variables on the balances are marginal resource values. Under convex technologies and competitive implementation they can support prices. For a distorted or distributionally weighted economy these dual variables do not automatically equal observed household prices or national welfare changes. Transfers from consumers to domestic producers must be counted once through ownership and household budgets.

\section*{Checks and unresolved part}
The bounds treat supplier concentration, correlation and shipping restrictions jointly; weak monotonicity can be strict or zero depending on redundant routes. Scarcity thresholds are in the same physical units as delivered end use. The exact calculation requires the joint structural shock law, while marginal outage frequencies yield only the demonstrated bounds. Investment timing can be incorporated, but investment incentives, substitution elasticities, informal rerouting, market clearing and welfare identification have not been solved. The complete dynamic general-equilibrium research target remains unresolved.

\section{Completion Estimate}
29\%

This is a post hoc, heuristic estimate for the full catalogue target.
The attempt derives physical scarcity bounds from time-expanded commodity networks, sharp supplier-dependence risk comparisons, and restrictions imposed by investment lags and linked production. Full equilibrium prices, household welfare, substitution, investment incentives, informal rerouting, and identification of dynamic fragmentation effects remain unresolved.

\section*{Source checked}
Jorge A. Alvarez, Mehdi Benatiya Andaloussi, Chiara Maggi, Alexandre Sollaci, Martin Stuermer and Petia Topalova (2023), \emph{Geoeconomic Fragmentation and Commodity Markets}, IMF Working Paper 23/201, Section 10. The accessible primary PDF's conclusion was inspected for its partial-equilibrium scope and extension agenda. \url{https://www.imf.org/-/media/files/publications/wp/2023/english/wpiea2023201-print-pdf.pdf}.

\end{document}
